\section{Additional Experimental Results}~\label{appx:exp}

The main text reports aggregate results for the threshold sweep and policy comparison. This appendix summarizes the design logic behind these aggregate values.

\paragraph{Threshold sweep.}
The threshold $\tau$ controls when belief confidence is considered sufficient for autonomous execution. Low thresholds allow the robot to commit to a state estimate with little feedback, which reduces query frequency and favors autonomy. Higher thresholds request more feedback and concentrate the belief before downstream actions are selected. This produces the saturation pattern in Table~\ref{tab:sys_tau}: both domains improve sharply around $\tau=0.7$, while larger values mainly increase query count after the task-relevant hypotheses have already been resolved.

\paragraph{Policy comparison.}
The policy comparison in Table~\ref{tab:system_eval} uses $\tau=0.8$ for \textbf{Ours}. \textbf{All} represents an upper-interaction baseline that asks whenever a query is available. \textbf{No} represents fully autonomous execution without corrective feedback. \textbf{Random} uses the same query interface but ignores belief confidence when deciding whether to ask. Thus, the comparison separates three effects: availability of human information, timing of feedback, and information gain of the selected hypothesis. The strong performance of \textbf{Ours} relative to \textbf{Random}, especially in tomato harvesting, shows the value of aligning feedback with the current belief ambiguity.

\paragraph{Domain difference.}
Waste sorting contains short pick-and-place sequences and requires correct category assignment before placement. Tomato harvesting contains longer action chains in which ambiguous tomato state estimates affect later scan, place, and discard decisions. Consequently, resolving belief errors at the right moment has greater leverage in tomato harvesting, while waste sorting gives shorter opportunities for uncertainty to propagate.
